\chapter{Construire un jeu de plateau, de A à Z}

Ce chapitre ne décrit pas un jeu : il le \textbf{construit}. Vous partez d'un
dossier vide et vous arrivez à un jeu de dames complet --- règles, damier,
adversaire, écran d'ouverture --- qui se lance sur Windows, Linux, macOS,
Android, iOS, HarmonyOS et dans un navigateur.

Treize étapes, dans l'ordre. Chacune produit quelque chose qui \emph{se
construit et se lance} : à aucun moment vous n'écrivez du code que vous ne
pouvez pas encore éprouver. Rien n'est sauté.

\begin{nkretenir}
\textbf{L'ordre des étapes est le sujet du chapitre.}

On écrit les règles \emph{avant} l'écran, et le banc \emph{avant} le dessin.
Ce n'est pas de la discipline : c'est ce qui rend le jeu prouvable. Écrivez le
dessin d'abord, et vos règles n'auront plus jamais de verdict indépendant de ce
que vous croyez voir.
\end{nkretenir}

\section*{Ce que nous allons construire}
\addcontentsline{toc}{section}{Ce que nous allons construire}

\begin{center}
\begin{tabular}{@{}rlll@{}}
\toprule
\textbf{Étape} & \textbf{Fichier} & \textbf{Ce qu'on obtient} & \textbf{Lignes} \\
\midrule
1 & \texttt{NkDames.jenga} & le projet existe & $\approx$60 \\
2 & \texttt{src/main.cpp} & une fenêtre qui s'ouvre & 12 \\
3--7 & \texttt{Dames/NkDamesRegles.\{h,cpp\}} & le jeu, sans image & 538 \\
8 & \texttt{Dames/NkDamesBanc.\{h,cpp\}} & \textbf{la preuve} & 272 \\
9 & \texttt{Dames/NkDamesTheme.h} & les couleurs & 63 \\
9--10 & \texttt{Dames/NkDamesEcran.\{h,cpp\}} & le damier à l'écran & 609 \\
11--13 & \texttt{Dames/NkDamesJeu.\{h,cpp\}} & le jeu jouable & 477 \\
\bottomrule
\end{tabular}
\end{center}

\noindent Total : \textbf{1\,998 lignes}, sept fichiers de code. Les chiffres
sont ceux de \texttt{Applications/NkDames} dans le dépôt --- ce chapitre en est
la reconstruction pas à pas.

\section{Étape 1 --- le dossier et le fichier de projet}

\begin{nkcode}{L'arborescence a creer}
Applications/NkDames/
    NkDames.jenga
    src/
        main.cpp
        Dames/
\end{nkcode}

Rien d'autre : pas de dossier \texttt{assets}, pas de \texttt{Resources}. Ce jeu
ne chargera \textbf{aucun fichier à l'exécution} --- le damier est dessiné par
géométrie et la police est embarquée dans le binaire.

\begin{nkretenir}
\textbf{C'est ce qui permet de démarrer à l'identique sur sept plateformes.} Un
seul asset à copier, et il faut une chaîne de copie par plateforme, une
permission de lecture sur Android, un téléchargement de plus avant la première
image sur le Web. Ne créez ce dossier que le jour où vous avez vraiment un
fichier à y mettre.
\end{nkretenir}

Voici le fichier de projet \textbf{minimal et complet} pour construire sur
ordinateur :

\begin{nkcode}{Applications/NkDames/NkDames.jenga -- version minimale, fonctionnelle}
from Jenga import *
from jengaconfig import *

with project("NkDames"):
    windowedapp()                 # pas de console noire derriere la fenetre
    language("C++")
    cppdialect("C++17")
    location(".")

    files(["src/**.cpp"])         # tout le .cpp sous src/, recursivement

    nkentseudependson(
        ["NKEvent", "NKWindow", "NKCanvas", "NKImage", "NKFont", "NKGlad",
         "NKLogger", "NKMath", "NKTime", "NKStream", "NKGui", "NKFileSystem",
         "NKSerialization", "NKContainers", "NKMemory", "NKCore",
         "NKPlatform", "NKThreading", "NKAudio"],
        extra_includes=["src", "%{NKGlad.location}/include",
                        "%{wks.location}/Externals"])

    objdir("%{wks.location}/Build/Obj/%{cfg.buildcfg}-%{cfg.system}/%{prj.name}")
    targetdir("%{wks.location}/Build/Bin/%{cfg.buildcfg}-%{cfg.system}/%{prj.name}")

    apppublisher("Rihen Universe")
    appversion("0.1.0")
    licensefile("../../LICENSE")

    with filter("system:Windows"):
        windowedapp()
        usetoolchain(TC_WINDOWS)
        defines(["WIN32_LEAN_AND_MEAN", "_UNICODE", "UNICODE"])
        links(["user32", "gdi32", "opengl32", "dwmapi", "shell32",
               "ole32", "uuid", "avrt", "winmm"])

    with filter("config:Debug"):
        defines(["_DEBUG", "DEBUG"]); optimize("Off");   symbols(True)
    with filter("config:Release"):
        defines(["NDEBUG"]);          optimize("Speed"); symbols(False)
\end{nkcode}

\begin{nkretenir}
\textbf{\texttt{extra\_includes=["src", \dots]} est la ligne qu'on oublie}, et
son absence donne \texttt{fatal error: 'Dames/NkDamesRegles.h' file not found}
--- une erreur qui accuse votre inclusion alors que c'est le \emph{projet} qui
ne sait pas où chercher.

Les dix-neuf modules listés ne sont pas décoratifs : chacun manquant produit une
erreur d'\emph{édition de liens}, pas de compilation --- donc à la toute fin,
avec un message qui nomme un symbole et jamais un module.
\end{nkretenir}

\begin{nkpiege}
\textbf{Un projet qui n'est pas enregistré dans l'espace de travail n'existe
pas.} Ajoutez, dans \texttt{Nkentseu.jenga} à la racine :

\begin{nkcode}{Nkentseu.jenga}
with include("Applications/NkDames/NkDames.jenga"):
    pass
\end{nkcode}

Sans cette ligne, \texttt{jenga build --target NkDames} répond que la cible
n'existe pas --- et l'on relit son \texttt{.jenga} vingt fois, qui est juste.
\end{nkpiege}

\subsection{Les six autres plateformes}

Chaque plateforme ajoute un \texttt{filter} qui déclare sa chaîne d'outils et
ses bibliothèques. Le tableau suivant résume ce que chacun apporte ; le fichier
complet est dans le dépôt.

\begin{center}
\begin{tabular}{@{}lll@{}}
\toprule
\textbf{Plateforme} & \textbf{Chaîne} & \textbf{À déclarer en plus} \\
\midrule
Linux (XLib) & \texttt{clang-native} & \texttt{X11 Xext Xrandr GL asound} \\
macOS & \texttt{clang-native} & \emph{frameworks} Cocoa, Metal, AudioToolbox\dots \\
Android & \texttt{android-ndk} & id, SDK, ABI, orientation, \texttt{OpenSLES} \\
iOS & --- & \emph{frameworks} UIKit, Metal, AudioToolbox \\
HarmonyOS & \texttt{ohos-ndk} & \texttt{bundlename}, \texttt{hilog\_ndk.z} \\
Web & \texttt{emscripten} & \texttt{ASYNCIFY}, \texttt{USE\_WEBGL2} \\
\bottomrule
\end{tabular}
\end{center}

\begin{nkpiege}[Trois erreurs de ce tableau ont réellement été payées]
\textbf{\texttt{AudioUnit} n'existe pas sur iOS.} C'est un framework séparé sur
macOS seulement ; sur iOS ses symboles vivent dans \texttt{AudioToolbox}.
Recopier la liste macOS donne \texttt{ld: framework AudioUnit not found}.
\emph{Un précédent se vérifie sur SA plateforme.}

\textbf{Sans orientation déclarée}, Android et HarmonyOS imposent la leur et
l'image sort tournée d'un quart de tour.

\textbf{\texttt{ASYNCIFY} n'est pas une option de confort} sur le Web : sans
elle, la boucle ne peut pas rendre la main et l'onglet gèle.
\end{nkpiege}

\section{Étape 2 --- la fenêtre qui s'ouvre}

Douze lignes, et le jeu se lance déjà.

\begin{nkcode}{Applications/NkDames/src/main.cpp}
#include "Dames/NkDamesJeu.h"
#include "NKWindow/NKMain.h"

NKENTSEU_DEFINE_APP_DATA(([]() {
    nkentseu::NkAppData d{};
    d.appName    = "NkDames";
    d.appVersion = "1.0.0";
    return d;
})());

int nkmain(const nkentseu::NkEntryState &state) {
    return nkentseu::renderer::NkCanvasApp::Run<nkentseu::jeux::dames::NkDamesJeu>(state);
}
\end{nkcode}

\begin{nkretenir}
\textbf{Ce fichier ne contiendra jamais de logique.} Il déclare le nom et la
version, puis rend la main à la coquille. Tout le reste a son fichier --- et
c'est cette discipline qui fait qu'à l'étape 13 vous saurez encore où ajouter
quelque chose.
\end{nkretenir}

\begin{nkpiege}
\textbf{\texttt{\#include "NKWindow/NKMain.h"} fournit le point d'entrée natif}
de chaque plateforme : \texttt{WinMain}, \texttt{android\_main},
\texttt{UIApplicationMain}\dots

Sans lui, tout compile et l'édition de liens tombe sur
\texttt{undefined reference to WinMain} --- une erreur qui ne nomme
\emph{aucun} de vos fichiers, et qu'on cherche donc partout sauf ici.

Et \textbf{n'écrivez pas votre propre \texttt{main}} : \texttt{nkmain} est le
point d'entrée portable, avec seize points d'entrée de plateforme derrière lui.
\end{nkpiege}

\begin{nkpiege}
\texttt{NKENTSEU\_DEFINE\_APP\_DATA} attend un \texttt{NkAppData}, pas une
chaîne. Lui passer \texttt{"NkDames"} donne
\texttt{no viable overloaded '='} --- un message qui ne nomme jamais la cause.

Le \texttt{()} final de la lambda n'est pas décoratif non plus : sans lui, vous
passez la \emph{fonction} au lieu de son résultat.
\end{nkpiege}

\section{Étape 3 --- les règles : les types}

Créez \texttt{src/Dames/NkDamesRegles.h}. Regardez d'abord ce qu'il inclut, car
c'est la décision la plus importante du chapitre :

\begin{nkcode}{Applications/NkDames/src/Dames/NkDamesRegles.h}
#include "NKContainers/Sequential/NkVector.h"
#include "NKCore/NkTypes.h"
\end{nkcode}

\begin{nkretenir}
\textbf{Deux inclusions. Ni fenêtre, ni GPU, ni image, ni police.}

Conséquence immédiate, et c'est tout l'objet de l'étape 8 : \textbf{ces règles se
testeront sans ouvrir de fenêtre}, en console, en une seconde, sur une machine
sans écran.

Rangez une seule règle du côté du dessin, et elle devient invisible au banc. Le
jeu marchera peut-être ; personne ne pourra le prouver.
\end{nkretenir}

\begin{nkcode}{Applications/NkDames/src/Dames/NkDamesRegles.h}
namespace nkentseu {
    namespace jeux {

        enum class NkDamesCamp : uint8 { NK_BLANC = 0, NK_NOIR = 1 };

        enum class NkDamesPiece : uint8 {
            NK_VIDE = 0,
            NK_PION_BLANC, NK_PION_NOIR,
            NK_DAME_BLANCHE, NK_DAME_NOIRE
        };

        inline bool NkDamesEstBlanc(NkDamesPiece p) noexcept {
            return p == NkDamesPiece::NK_PION_BLANC || p == NkDamesPiece::NK_DAME_BLANCHE;
        }
        inline bool NkDamesEstNoir(NkDamesPiece p) noexcept {
            return p == NkDamesPiece::NK_PION_NOIR || p == NkDamesPiece::NK_DAME_NOIRE;
        }
        inline bool NkDamesEstDame(NkDamesPiece p) noexcept {
            return p == NkDamesPiece::NK_DAME_BLANCHE || p == NkDamesPiece::NK_DAME_NOIRE;
        }
        inline bool NkDamesAppartient(NkDamesPiece p, NkDamesCamp c) noexcept {
            return c == NkDamesCamp::NK_BLANC ? NkDamesEstBlanc(p) : NkDamesEstNoir(p);
        }
\end{nkcode}

\begin{nkretenir}
\textbf{Une seule énumération pour la pièce, pas deux champs} « couleur » et
« type ». Deux champs autorisent l'état \emph{aucune couleur mais une dame} ---
un état qui n'existe pas dans le jeu et qu'il faudrait pourtant traiter partout.

Les quatre fonctions \texttt{inline} sont la contrepartie : elles rendent
lisible ce que l'énumération compacte. Elles sont dans l'en-tête, donc
gratuites.
\end{nkretenir}

Puis le coup. Un coup de dames n'est pas un déplacement : c'est \emph{une rafle
entière}.

\begin{nkcode}{Applications/NkDames/src/Dames/NkDamesRegles.h}
        static const int32 NK_DAMES_MAX_PRISES = 20;

        struct NkDamesCoup {
                int8 depR = -1, depC = -1;   // case de depart
                int8 arrR = -1, arrC = -1;   // case d'arrivee FINALE

                uint8 nbPrises = 0;
                int8 prisR[NK_DAMES_MAX_PRISES] = {};
                int8 prisC[NK_DAMES_MAX_PRISES] = {};

                // Les cases d'atterrissage successives, pour l'affichage.
                uint8 nbEtapes = 0;
                int8 etapeR[NK_DAMES_MAX_PRISES] = {};
                int8 etapeC[NK_DAMES_MAX_PRISES] = {};

                bool EstPrise() const noexcept { return nbPrises > 0; }
        };
\end{nkcode}

\begin{nkretenir}
\textbf{Des tableaux de taille fixe, pas des vecteurs.} Un coup est copié des
milliers de fois pendant la génération ; une allocation par copie coûterait
plus que tout le reste du jeu.

Le commentaire du dépôt justifie le 20 : « le record théorique connu au jeu
international est de 12 ; 20 laisse une marge sans jamais tronquer une rafle
légale --- et une troncature silencieuse fausserait le choix de la rafle
maximale, donc la règle elle-même ».

\textbf{C'est un plafond qui se justifie par la conséquence de son
dépassement}, pas par une intuition de confort.
\end{nkretenir}

Enfin la classe :

\begin{nkcode}{Applications/NkDames/src/Dames/NkDamesRegles.h}
        class NkDamesPartie {
            public:
                static const int32 NK_TAILLE = 10;

                NkDamesPartie() { Initialiser(); }

                void Initialiser() noexcept;

                NkDamesPiece Case(int32 r, int32 c) const noexcept {
                    return DansDamier(r, c) ? mCases[r][c] : NkDamesPiece::NK_VIDE;
                }
                void PoserCase(int32 r, int32 c, NkDamesPiece p) noexcept {
                    if (DansDamier(r, c)) { mCases[r][c] = p; }
                }

                NkDamesCamp Trait() const noexcept { return mTrait; }
                void PoserTrait(NkDamesCamp c) noexcept { mTrait = c; }

                void CoupsLegaux(NkVector<NkDamesCoup> &sortie) const;
                bool Jouer(const NkDamesCoup &coup) noexcept;
                bool EstTerminee(NkDamesCamp &gagnant) const;
                int32 CompterPieces(NkDamesCamp camp) const noexcept;
                void CoupsDepuis(int32 r, int32 c, NkVector<NkDamesCoup> &sortie) const;

                static bool DansDamier(int32 r, int32 c) noexcept {
                    return r >= 0 && r < NK_TAILLE && c >= 0 && c < NK_TAILLE;
                }
                // Les cases jouables sont les SOMBRES. Convention : (r+c) impair.
                static bool CaseJouable(int32 r, int32 c) noexcept {
                    return ((r + c) & 1) != 0;
                }

            private:
                void GenererPrises(int32 r, int32 c, NkVector<NkDamesCoup> &sortie) const;
                void ExplorerPrises(int32 r, int32 c, NkDamesPiece piece, NkDamesCoup &courant,
                                    bool prisesFaites[NK_TAILLE][NK_TAILLE],
                                    NkVector<NkDamesCoup> &sortie) const;
                void GenererSimples(int32 r, int32 c, NkVector<NkDamesCoup> &sortie) const;

                NkDamesPiece mCases[NK_TAILLE][NK_TAILLE] = {};
                NkDamesCamp mTrait = NkDamesCamp::NK_BLANC;
        };
\end{nkcode}

\begin{nkretenir}
\textbf{Six questions publiques suffisent à décrire un jeu au tour par tour} :
initialiser, lire une case, qui joue, quels coups sont légaux, jouer, est-ce
fini. Vous retrouverez exactement les mêmes aux échecs et au ludo.

\texttt{Case()} borne elle-même : un indice hors damier rend \texttt{NK\_VIDE}
au lieu de lire n'importe où. C'est ce qui permet aux générateurs de coups
d'explorer sans tester les bords à chaque pas.

Et \texttt{CoupsLegaux} \emph{écrit dans} un vecteur qu'on lui donne au lieu
d'en rendre un : l'appelant réutilise le même tampon à chaque tour.
\end{nkretenir}

\section{Étape 4 --- l'initialisation}

Créez \texttt{NkDamesRegles.cpp}. Première fonction :

\begin{nkcode}{Applications/NkDames/src/Dames/NkDamesRegles.cpp}
void NkDamesPartie::Initialiser() noexcept {
    for (int32 r = 0; r < NK_TAILLE; ++r) {
        for (int32 c = 0; c < NK_TAILLE; ++c) {
            mCases[r][c] = NkDamesPiece::NK_VIDE;
        }
    }
    // Rangees 0..3 = noirs (en haut), 6..9 = blancs (en bas).
    for (int32 r = 0; r < 4; ++r) {
        for (int32 c = 0; c < NK_TAILLE; ++c) {
            if (CaseJouable(r, c)) { mCases[r][c] = NkDamesPiece::NK_PION_NOIR; }
        }
    }
    for (int32 r = 6; r < NK_TAILLE; ++r) {
        for (int32 c = 0; c < NK_TAILLE; ++c) {
            if (CaseJouable(r, c)) { mCases[r][c] = NkDamesPiece::NK_PION_BLANC; }
        }
    }
    mTrait = NkDamesCamp::NK_BLANC;
}
\end{nkcode}

\begin{nkretenir}
\textbf{Les rangées 4 et 5 restent vides} : c'est ce qui donne 20 pièces par
camp sur un damier de 10$\times$10. Et \texttt{CaseJouable} est appelée plutôt
que réécrite --- une seconde définition de « case sombre » quelque part dans le
dessin, et un jour les deux diffèrent.
\end{nkretenir}

Et le compteur, qui servira à la fois aux règles, à l'affichage et au banc :

\begin{nkcode}{Applications/NkDames/src/Dames/NkDamesRegles.cpp}
int32 NkDamesPartie::CompterPieces(NkDamesCamp camp) const noexcept {
    int32 n = 0;
    for (int32 r = 0; r < NK_TAILLE; ++r) {
        for (int32 c = 0; c < NK_TAILLE; ++c) {
            if (NkDamesAppartient(mCases[r][c], camp)) { ++n; }
        }
    }
    return n;
}
\end{nkcode}

\section{Étape 5 --- les coups simples}

En haut du \texttt{.cpp}, les quatre diagonales :

\begin{nkcode}{Applications/NkDames/src/Dames/NkDamesRegles.cpp}
namespace {
    // Ordre fixe : il decide, a rafle egale, quel coup l'ordinateur trouve
    // en premier -- donc il doit etre stable.
    const int32 kDirR[4] = {-1, -1, 1, 1};
    const int32 kDirC[4] = {-1, 1, -1, 1};
}
\end{nkcode}

\begin{nkretenir}
\textbf{Un ordre arbitraire qui a une conséquence observable cesse d'être
arbitraire.} Changer ces quatre paires ne change aucune règle --- mais change le
coup que l'IA choisit à égalité, donc la partie entière.

Une constante dont la valeur n'importe pas mais dont \emph{la stabilité}
importe : écrivez-le, sinon quelqu'un les réordonnera pour faire joli.
\end{nkretenir}

Un pion avance d'une case en diagonale, sans reculer. Une dame glisse.

\begin{nkcode}{Applications/NkDames/src/Dames/NkDamesRegles.cpp -- structure}
void NkDamesPartie::GenererSimples(int32 r, int32 c, NkVector<NkDamesCoup> &sortie) const {
    const NkDamesPiece piece = mCases[r][c];
    const bool dame  = NkDamesEstDame(piece);
    const int32 sens = NkDamesEstBlanc(piece) ? -1 : +1;   // les blancs montent

    for (int32 d = 0; d < 4; ++d) {
        if (!dame && kDirR[d] != sens) { continue; }       // un pion ne recule pas
        int32 rr = r + kDirR[d], cc = c + kDirC[d];
        while (DansDamier(rr, cc) && mCases[rr][cc] == NkDamesPiece::NK_VIDE) {
            NkDamesCoup m;
            m.depR = (int8)r;  m.depC = (int8)c;
            m.arrR = (int8)rr; m.arrC = (int8)cc;
            sortie.PushBack(m);
            if (!dame) { break; }                          // le pion s'arrete a 1
            rr += kDirR[d]; cc += kDirC[d];                // la dame glisse
        }
    }
}
\end{nkcode}

\begin{nkretenir}
\textbf{Une seule fonction pour le pion et la dame}, séparés par deux lignes :
le \texttt{continue} qui interdit le recul, et le \texttt{break} qui limite à une
case. Deux fonctions distinctes auraient divergé au premier correctif.
\end{nkretenir}

\section{Étape 6 --- la rafle, et c'est le seul endroit délicat}

Une rafle est une \emph{suite} de prises dans le même coup. On l'explore en
profondeur.

\begin{nkcode}{Applications/NkDames/src/Dames/NkDamesRegles.cpp -- l'en-tete du fichier}
// LE POINT DELICAT, ET IL N'Y EN A QU'UN : la rafle.
//   On l'explore en profondeur, et la piece prise est marquee "deja prise"
//   pendant l'exploration SANS ETRE RETIREE du damier : elle continue de
//   BLOQUER le passage. C'est la regle du jeu, et c'est aussi ce qui evite de
//   compter deux fois la meme piece. Retirer la piece pendant l'exploration
//   produirait des rafles qui n'existent pas -- l'erreur classique, et elle
//   est silencieuse.
\end{nkcode}

\begin{nkretenir}
\textbf{Lisez ce paragraphe deux fois : il contient la totalité de la
difficulté.}

Une pièce capturée n'est retirée qu'\emph{à la fin} de la rafle. Pendant
l'exploration elle est marquée --- pour ne pas être comptée deux fois --- mais
elle reste sur le damier et \textbf{bloque le trajet}.

Si vous la retirez tout de suite, votre programme trouvera des rafles qui
n'existent pas : la case libérée laisse passer une dame qui, en vrai, serait
arrêtée. Le jeu ne plantera pas. Il trichera.
\end{nkretenir}

\begin{nkcode}{Applications/NkDames/src/Dames/NkDamesRegles.cpp}
void NkDamesPartie::ExplorerPrises(int32 r, int32 c, NkDamesPiece piece,
                                   NkDamesCoup &courant,
                                   bool prisesFaites[NK_TAILLE][NK_TAILLE],
                                   NkVector<NkDamesCoup> &sortie) const {
    const NkDamesCamp camp = NkDamesEstBlanc(piece) ? NkDamesCamp::NK_BLANC
                                                   : NkDamesCamp::NK_NOIR;
    const NkDamesCamp adverse = (camp == NkDamesCamp::NK_BLANC) ? NkDamesCamp::NK_NOIR
                                                                : NkDamesCamp::NK_BLANC;
    const bool dame = NkDamesEstDame(piece);
    bool aContinue = false;

    for (int32 d = 0; d < 4; ++d) {
        // 1. avancer jusqu'a rencontrer une piece
        // 2. si elle est adverse ET pas deja prise ET la case suivante est libre :
        //      -> marquer, empiler l'etape, RECURSER, puis DEMARQUER
        //      -> aContinue = true
    }

    // 3. une rafle qui ne peut plus continuer est un coup complet
    if (!aContinue && courant.nbPrises > 0) {
        courant.arrR = (int8)r; courant.arrC = (int8)c;
        sortie.PushBack(courant);
    }
}
\end{nkcode}

\begin{nkretenir}
\textbf{Le « démarquer » après la récursion est ce qui distingue une exploration
d'un parcours.} Sans lui, la première branche essayée empoisonne toutes les
suivantes : la deuxième croit que la pièce est déjà prise, et la rafle qu'elle
aurait trouvée disparaît.

C'est le \emph{retour sur trace} (\emph{backtracking}), et il tient en une
ligne qu'on oublie une fois sur deux.
\end{nkretenir}

\begin{nkpiege}
\textbf{La case de DÉPART de la rafle est vide « pour de faux ».} La pièce y est
encore, physiquement, puisqu'on n'a rien joué. Une dame qui revient sur sa
propre case de départ en cours de rafle doit pouvoir passer.

C'est le genre de cas qu'on ne rencontre qu'une partie sur cent --- donc jamais
pendant le développement, et toujours chez un joueur.
\end{nkpiege}

\section{Étape 7 --- rafle maximale, jouer, fin de partie}

\begin{nkcode}{Applications/NkDames/src/Dames/NkDamesRegles.cpp -- structure}
void NkDamesPartie::CoupsLegaux(NkVector<NkDamesCoup> &sortie) const {
    sortie.Clear();
    // 1. generer TOUTES les prises de toutes les pieces du camp au trait
    // 2. s'il y en a : garder UNIQUEMENT celles de nbPrises maximal
    // 3. sinon seulement : generer les coups simples
}
\end{nkcode}

\begin{nkretenir}
\textbf{Le filtrage par maximum se fait EN SORTIE, jamais pendant
l'exploration.} Filtrer trop tôt fait perdre des branches qui allaient devenir
plus longues : une rafle de 2 prises peut se prolonger en 5.

C'est une erreur d'optimisation prématurée qui donne un jeu \emph{presque}
juste --- le pire cas, car il faut une position précise pour la voir.
\end{nkretenir}

\begin{nkcode}{Applications/NkDames/src/Dames/NkDamesRegles.cpp -- structure}
bool NkDamesPartie::Jouer(const NkDamesCoup &coup) noexcept {
    NkDamesPiece p = Case(coup.depR, coup.depC);
    if (!NkDamesAppartient(p, mTrait)) {
        return false;                       // un refus se DIT
    }
    PoserCase(coup.depR, coup.depC, NkDamesPiece::NK_VIDE);
    for (uint8 i = 0; i < coup.nbPrises; ++i) {
        PoserCase(coup.prisR[i], coup.prisC[i], NkDamesPiece::NK_VIDE);
    }
    // Promotion : SEULEMENT si le coup S'ACHEVE sur la derniere rangee.
    if (!NkDamesEstDame(p)) {
        if (NkDamesEstBlanc(p) && coup.arrR == 0)              { p = NkDamesPiece::NK_DAME_BLANCHE; }
        else if (NkDamesEstNoir(p) && coup.arrR == NK_TAILLE-1) { p = NkDamesPiece::NK_DAME_NOIRE; }
    }
    PoserCase(coup.arrR, coup.arrC, p);
    mTrait = (mTrait == NkDamesCamp::NK_BLANC) ? NkDamesCamp::NK_NOIR : NkDamesCamp::NK_BLANC;
    return true;
}
\end{nkcode}

\begin{nkretenir}
\textbf{\texttt{Jouer} rend un booléen, et il faut le lire.} « Un refus se dit,
il ne se devine pas à l'absence d'effet » : sans ce retour, un coup illégal
laisse le damier inchangé et le joueur croit que son clic n'a pas été reçu.

\textbf{Et la promotion n'a lieu que si le coup S'ACHÈVE sur la dernière
rangée.} Un pion qui traverse cette rangée \emph{au milieu} d'une rafle ne
devient pas dame --- règle réelle, et l'une des cinq que toutes les
implémentations ratent.
\end{nkretenir}

\begin{nkcode}{Applications/NkDames/src/Dames/NkDamesRegles.cpp -- structure}
bool NkDamesPartie::EstTerminee(NkDamesCamp &gagnant) const {
    if (CompterPieces(NkDamesCamp::NK_BLANC) == 0) { gagnant = NkDamesCamp::NK_NOIR;  return true; }
    if (CompterPieces(NkDamesCamp::NK_NOIR)  == 0) { gagnant = NkDamesCamp::NK_BLANC; return true; }
    NkVector<NkDamesCoup> coups;
    CoupsLegaux(coups);
    if (coups.Size() == 0) {          // bloque = perdu
        gagnant = (mTrait == NkDamesCamp::NK_BLANC) ? NkDamesCamp::NK_NOIR : NkDamesCamp::NK_BLANC;
        return true;
    }
    return false;
}
\end{nkcode}

\begin{nkpiege}
\textbf{Un camp sans coup légal a perdu, même s'il lui reste des pièces.} Oublier
ce cas donne un jeu qui \emph{se fige} : c'est au joueur de jouer, il ne peut
rien jouer, et rien ne le dit. Le symptôme --- « le jeu ne répond plus » --- ne
ressemble en rien à sa cause.
\end{nkpiege}

\section{Étape 8 --- le banc. Avant l'écran.}

\textbf{Vous n'avez encore dessiné aucun pixel, et vous allez déjà prouver que
votre jeu est juste.} C'est possible uniquement parce que l'étape 3 n'a inclus
ni fenêtre ni GPU.

\begin{nkcode}{Applications/NkDames/src/Dames/NkDamesBanc.h}
namespace nkentseu { namespace jeux { namespace dames {
    /// Rend 0 si tout passe, non nul sinon. Le code de sortie EST le verdict.
    int32 NkDamesLancerBanc();
}}}
\end{nkcode}

Et on le branche \emph{avant} toute fenêtre, dans le jeu :

\begin{nkcode}{Applications/NkDames/src/Dames/NkDamesJeu.cpp -- structure}
NkOptional<int> NkDamesJeu::OnCommandLine(const NkVector<NkString> &args) {
    for (uint64 i = 0; i < args.Size(); ++i) {
        if (args[i] == "--selftest") {
            return NkDamesLancerBanc();     // on rend un code : pas de fenetre
        }
    }
    return {};                              // rien rendu = on continue normalement
}
\end{nkcode}

\begin{nkretenir}
\textbf{\texttt{OnCommandLine} s'exécute AVANT la création de la fenêtre.} C'est
ce qui permet à \texttt{NkDames --selftest} de tourner dans une intégration
continue sans écran, sans GPU, en une seconde.

Rendre une valeur arrête le programme avec ce code ; ne rien rendre le laisse
démarrer. Un \texttt{NkOptional} plutôt qu'un \texttt{int} sentinelle : « pas de
réponse » et « le code 0 » sont deux choses différentes.
\end{nkretenir}

Un cas de banc ressemble à ceci --- on \emph{fabrique} la position :

\begin{nkcode}{Applications/NkDames/src/Dames/NkDamesBanc.cpp -- forme d'un cas}
// Cas : la promotion n'a lieu QUE si le coup s'acheve sur la derniere rangee.
{
    NkDamesPartie p;
    for (int32 r = 0; r < 10; ++r)
        for (int32 c = 0; c < 10; ++c)
            p.PoserCase(r, c, NkDamesPiece::NK_VIDE);   // damier vide : on maitrise tout
    p.PoserCase(2, 1, NkDamesPiece::NK_PION_BLANC);
    p.PoserCase(1, 2, NkDamesPiece::NK_PION_NOIR);      // a sauter, arrivee rangee 0
    p.PoserCase(1, 4, NkDamesPiece::NK_PION_NOIR);      // la rafle CONTINUE
    p.PoserTrait(NkDamesCamp::NK_BLANC);

    NkVector<NkDamesCoup> coups;
    p.CoupsLegaux(coups);
    VERIFIER(coups.Size() == 1);
    p.Jouer(coups[0]);
    VERIFIER(p.Case(coups[0].arrR, coups[0].arrC) == NkDamesPiece::NK_PION_BLANC);
    //        ^ un PION, pas une dame : la rafle ne s'acheve pas rangee 0
}
\end{nkcode}

\begin{nkretenir}
\textbf{On part d'un damier VIDE, jamais de la position initiale.} Sur un damier
complet, vingt autres coups sont légaux et l'on ne sait plus lequel on mesure.

Fabriquer la position est ce qui rend l'assertion \emph{lisible} :
\texttt{coups.Size() == 1} n'a de sens que parce qu'on a construit la situation
où il n'y a qu'un coup.
\end{nkretenir}

\begin{nkpiege}[Les cas négatifs comptent autant que les positifs]
Un banc qui vérifie seulement que les coups légaux sont acceptés passe au vert
sur des règles qui acceptent \textbf{tout}.

Écrivez donc aussi : \emph{on ne peut pas jouer une pièce adverse},
\emph{une rafle de 3 est obligatoire quand une de 2 existe},
\emph{un pion ne recule pas en déplacement simple}.
\end{nkpiege}

\begin{nkretenir}[La contre-épreuve, et elle n'est pas facultative]
\textbf{Un banc qui n'a jamais rougi est une intention, pas un contrôle.}

Cassez exprès ce qu'il surveille --- retirez le filtre de rafle maximale --- et
vérifiez qu'il passe au \textbf{rouge}, sur le bon cas, avec un code de sortie
non nul. Puis restaurez.

Le dépôt inscrit cette contre-épreuve dans l'en-tête du banc, avec le cas qui
est tombé :

\begin{nkcode}{Applications/NkDames/src/Dames/NkDamesBanc.cpp}
// CONTRE-EPREUVE FAITE : en retirant le filtre de rafle maximale, le banc
// est passe au ROUGE sur le bon cas, code de sortie 1.
\end{nkcode}
\end{nkretenir}

\begin{nkpiege}[Le cas qui a été faux avant d'être juste]
Le banc du ludo exigeait d'abord que chaque pas de la piste soit
\emph{orthogonal}. Il échouait --- et le code avait raison : un vrai plateau de
ludo a exactement quatre virages en diagonale, aux coins.

\textbf{Un banc rouge n'accuse pas toujours le code.} Avant de corriger ce qui
est mesuré, vérifiez ce que la mesure attendait.
\end{nkpiege}

\noindent Lancez-le : \texttt{NkDames.exe --selftest}. Vous avez un jeu prouvé,
et toujours aucun pixel.

\section{Étape 9 --- le thème, puis la géométrie}

Une seule source pour les couleurs, dans un en-tête sans code :

\begin{nkcode}{Applications/NkDames/src/Dames/NkDamesTheme.h -- forme}
struct NkDamesTheme {
        NkColor fond, caseClaire, caseSombre, liseréSel, dispo;
        NkColor pionBlanc, pionNoir, contour, texte;
};
inline const NkDamesTheme &NkDamesThemeDefaut();
\end{nkcode}

\begin{nkretenir}
\textbf{Aucune couleur écrite ailleurs.} Une seule valeur en dur dans le dessin,
et le jour où l'on veut un thème clair il faut la retrouver --- ce qui se
termine toujours par « il en restait une ».
\end{nkretenir}

Puis la géométrie. \textbf{Elle se calcule, elle ne se code pas en dur} :

\begin{nkcode}{Applications/NkDames/src/Dames/NkDamesEcran.h -- forme}
struct NkDamesGeometrie {
        NkRect damier;      // le carre du jeu
        float32 caseTaille; // cote d'une case
        NkRect hud;         // la bande d'information
        NkRect boutons[2];  // les sieges

        static NkDamesGeometrie Calculer(const renderer::NkLayoutInfo &info);
};
\end{nkcode}

\begin{nkpiege}
\textbf{Le bloc se centre sur l'ÉCRAN, pas sur le plateau.} Le ludo a d'abord
ancré son menu sur le damier --- un carré centré --- et laissait une bande vide
sur tout le tiers supérieur en fenêtre portrait.

Chaque rectangle était juste ; leur \emph{ensemble} était mal placé. C'est un
défaut qu'aucun calcul ne signale et qui ne se voit que sur une image.
\end{nkpiege}

\begin{nkretenir}
\texttt{NkLayoutInfo} porte la \textbf{zone sûre} --- les marges de l'encoche et
de l'indicateur de geste. Partez d'elle, jamais de la taille brute de la fenêtre,
sinon la première rangée du damier passe sous l'encoche sur téléphone.
\end{nkretenir}

\section{Étape 10 --- le dessin}

L'écran reçoit une \emph{vue} : ce qu'il doit montrer, jamais l'objet du jeu.

\begin{nkcode}{Applications/NkDames/src/Dames/NkDamesEcran.h -- forme}
struct NkDamesVue {
        const NkDamesPartie *partie;
        const NkVector<NkDamesCoup> *coupsProposes;
        int32 selR, selC;
        bool finie;
        NkDamesCamp gagnant;
};

void NkDamesDessiner(nkgui::NkGuiDrawList &dl, const NkDamesGeometrie &geo,
                     const NkDamesVue &vue);
\end{nkcode}

\begin{nkretenir}
\textbf{Le dessin ne prend que des \texttt{const}.} Il ne peut donc rien
modifier --- ce n'est plus une discipline, c'est le compilateur qui l'impose.

Une fonction de dessin qui recevrait \texttt{NkDamesJeu\&} finirait, un jour de
fatigue, par jouer un coup depuis le dessin. Et le défaut serait introuvable :
personne ne cherche une règle de jeu dans un fichier d'affichage.
\end{nkretenir}

\begin{nkcode}{Applications/NkDames/src/Dames/NkDamesEcran.cpp -- ordre du dessin}
// 1. le fond
// 2. les 100 cases (claire / sombre)
// 3. le liseré de la case selectionnee
// 4. les pastilles des destinations possibles
// 5. les pieces (disque + anneau ; une couronne si c'est une dame)
// 6. le HUD : compteurs, trait, boutons de siege
// 7. le bandeau de fin, s'il y a lieu
\end{nkcode}

\begin{nkretenir}
\textbf{L'ordre EST le résultat.} Il n'y a pas de profondeur en 2D : ce qui est
dessiné en dernier est devant. Le liseré avant les pièces, et il disparaît
dessous.
\end{nkretenir}

\section{Étape 11 --- le clic devient un coup}

\begin{nkcode}{Applications/NkDames/src/Dames/NkDamesJeu.cpp -- structure}
bool NkDamesJeu::OnPointer(const NkPointer &p) {
    if (mSplash.Actif())            { return mSplash.Sauter(p); }
    if (p.phase != NkPointerPhase::NK_PRESSED) { return false; }

    // 1. le point est-il dans le damier ? -> (r, c)
    // 2. si une piece est deja selectionnee et (r,c) est une destination
    //       proposee -> jouer ce coup, deselectionner
    // 3. sinon, si (r,c) porte une piece du camp au trait
    //       -> selectionner, et demander CoupsDepuis(r, c, mCoupsProposes)
    // 4. sinon -> deselectionner
    return true;
}
\end{nkcode}

\begin{nkretenir}
\textbf{Deux clics, jamais un glisser.} Sélectionner puis poser marche à la
souris \emph{et} au doigt, sans code différent. Un glisser-déposer demanderait un
suivi de mouvement, un seuil de déclenchement, et une gestion de l'annulation.

\textbf{Et c'est \texttt{CoupsDepuis} qui alimente l'affichage des destinations}
--- la même fonction de règles qui décide. Recalculer les destinations dans
l'écran donnerait deux réponses qui divergeraient un jour.
\end{nkretenir}

\begin{nkpiege}
\texttt{NkPointer} donne des \textbf{pixels client}, pas des pixels écran. Et la
géométrie du damier est en pixels client aussi. Les mélanger donne une
sélection décalée d'exactement la position de la fenêtre --- donc juste quand la
fenêtre est en haut à gauche, et fausse partout ailleurs.
\end{nkpiege}

\section{Étape 12 --- l'adversaire}

\begin{nkcode}{Applications/NkDames/src/Dames/NkDamesRegles.h}
/// Choisit un coup pour l'ordinateur. Simple par construction : elle prefere
/// la rafle la plus grosse, puis la promotion, puis l'avance.
/// Ce n'est pas un adversaire fort, et le fichier ne pretend pas l'etre.
const NkDamesCoup *NkDamesChoisirCoup(const NkDamesPartie &partie,
                                      const NkVector<NkDamesCoup> &coups,
                                      uint32 &graine) noexcept;
\end{nkcode}

\begin{nkretenir}
\textbf{« Ce n'est pas un adversaire fort, et le fichier ne prétend pas
l'être. »} Une IA de trente lignes qui s'annonce comme telle est honnête ; la
même sans cette phrase déçoit.

Notez la \texttt{graine} passée \emph{par référence} : l'IA départage les
ex \ae quo au hasard, et ce hasard est \textbf{reproductible}. Même graine, même
partie --- donc un défaut d'IA se rejoue.
\end{nkretenir}

\begin{nkcode}{Applications/NkDames/src/Dames/NkDamesJeu.cpp -- structure}
void NkDamesJeu::OnTick(float32 deltaTime) {
    if (mSplash.Avancer(deltaTime)) { return; }   // l'ouverture tient l'ecran
    if (mAnim.Actif())              { mAnim.Avancer(deltaTime); return; }
    if (mFinie || EstHumain(TraitIndex())) { return; }

    mAttenteIA -= deltaTime;                       // laisser le temps de VOIR
    if (mAttenteIA <= 0.f) { JouerCoupIA(); mAttenteIA = 0.45f; }
}
\end{nkcode}

\begin{nkpiege}
\textbf{\texttt{Avancer} rend « vrai, je tiens l'écran », et il faut s'arrêter
là.} Sans ce retour, le jeu tournerait \emph{derrière} l'écran d'ouverture :
l'IA jouerait ses premiers coups pendant les quatre secondes, et le joueur
trouverait la partie déjà entamée en arrivant.

\textbf{Et une IA qui joue instantanément donne l'impression d'un bug.} Les
0,45 seconde d'attente ne servent pas l'IA : elles servent l'\oe il.
\end{nkpiege}

\section{Étape 13 --- l'ouverture, et les modes}

\begin{nkcode}{Applications/NkDames/src/Dames/NkDamesJeu.cpp}
bool NkDamesJeu::OnGuiInit() {
    mSplash.PoserJeu("Dames");
    Rejouer();
    return true;
}
\end{nkcode}

\begin{nkretenir}
\textbf{Une ligne pour l'écran d'ouverture.} La coquille porte la marque RIHEN
et l'animation ; vous ne donnez que le nom du jeu. C'est ce qui rend les
ouvertures des quatre jeux du dépôt identiques, à un mot près.

Et elle est \textbf{toujours sautable} : « un écran d'ouverture qu'on ne peut pas
passer est une taxe payée à chaque lancement, y compris par celui qui teste
vingt fois par heure ».
\end{nkretenir}

Enfin, les modes. Chaque siège a un contrôleur, et rien d'autre n'en découle :

\begin{nkcode}{Applications/NkDames/src/Dames/NkDamesJeu.h}
/// Par defaut : vous les blancs, l'IA les noirs.
NkControleur mControleur[2] = {NkControleur::NK_HUMAIN, NkControleur::NK_IA};
\end{nkcode}

\begin{nkretenir}
\textbf{Trois modes de jeu, et aucun n'est un cas particulier du code} :
humain/IA, humain/humain à deux sur le même écran, IA/IA en simulation. Un
tableau de deux contrôleurs les produit tous les trois.

Écrire \texttt{if (modeDuo) \dots} aurait donné trois chemins à tester au lieu
d'un.
\end{nkretenir}

\section{Construire, lancer, empaqueter}

\begin{nkcode}{Les commandes, dans l'ordre ou vous les taperez}
jenga build --target NkDames --config Release
.\Build\Bin\Release-Windows\NkDames\NkDames.exe --selftest    # le banc, 1 s
.\Build\Bin\Release-Windows\NkDames\NkDames.exe               # le jeu

jenga build --target NkDames --platform web
jenga package --platform windows --project NkDames --output dist/
\end{nkcode}

\begin{nkpiege}
\textbf{\texttt{jenga package} nomme sa sortie d'après le PROJET, pas d'après la
plateforme.} Empaqueter Windows puis Web dans le même dossier produit deux fois
\texttt{NkDames.zip}, et le second écrase le premier \emph{en silence}, code de
sortie 0.

Donnez un dossier de sortie par plateforme. Et \textbf{ouvrez l'archive} : c'est
en y trouvant un \texttt{.wasm} là où l'on attendait un \texttt{.exe} que le
défaut se voit, jamais dans la sortie de la commande.
\end{nkpiege}

\begin{nkretenir}
\textbf{Un APK Android de ce jeu n'a légitimement pas de \texttt{classes.dex}.}
C'est une \emph{NativeActivity} : il n'y a pas de code Java à convertir. Un
contrôle qui exige un \texttt{dex} déclarera cassé un APK qui s'installe et
fonctionne.
\end{nkretenir}

\section{Ce qui change pour les échecs et le ludo}

Le squelette ne bouge pas. Ce qui change tient dans une colonne.

\begin{center}
\begin{tabular}{@{}lll@{}}
\toprule
& \textbf{Ce qui est identique} & \textbf{Sa difficulté propre} \\
\midrule
Dames & les 7 fichiers, les 6 questions & la \emph{rafle maximale} \\
Échecs & les 7 fichiers, les 6 questions & roque, prise en passant, échec et mat \\
Ludo & les 7 fichiers, les 6 questions & 4 sièges, un dé, des sièges éteints \\
\bottomrule
\end{tabular}
\end{center}

\begin{center}
\begin{tabular}{@{}llll@{}}
\toprule
& \textbf{NkDames} & \textbf{NkEchecs} & \textbf{NkLudo} \\
\midrule
Fichiers & 10 & 9 & 9 \\
Lignes & 1\,998 & 2\,121 & 2\,130 \\
Cas de banc & 16 & 14 & 23 \\
\bottomrule
\end{tabular}
\end{center}

\begin{nkretenir}
\textbf{Trois jeux, trois cents lignes d'écart.} Ce n'est pas une coïncidence :
c'est le squelette qui domine, et la règle propre qui est petite.

Apprendre à construire l'un, c'est savoir construire les trois --- et c'est
exactement ce que ce chapitre voulait vous laisser.
\end{nkretenir}

\begin{nkretenir}[La leçon du ludo, qui dépasse le ludo]
Un siège peut être Humain, IA, ou \textbf{désactivé}. Sauter un siège éteint est
une règle de \emph{tour} --- elle vit donc dans les règles, pas dans l'écran.

Mise du côté du dessin, elle aurait été invisible au banc. Et le banc est le seul
endroit qui vérifie qu'une partie \textbf{se termine} : sans lui, un tour qui
donne la main à un siège sans pions ferait attendre le jeu indéfiniment --- et le
symptôme serait « le jeu se fige », à des lieues de sa cause.
\end{nkretenir}

\begin{nkbilan}
Treize étapes, sept fichiers, et un ordre qui n'est pas négociable : le projet,
la fenêtre, les règles, \textbf{le banc}, puis seulement l'écran, l'entrée,
l'adversaire. Les règles n'incluent que deux en-têtes --- c'est ce qui rend le
banc possible, et le banc est ce qui distingue un jeu qui marche d'un jeu dont
on \emph{sait} qu'il marche. Le dessin ne reçoit que du \texttt{const}, l'écran
ne décide de rien, et le fichier de jeu est le seul à modifier l'état. Les
échecs et le ludo se construisent avec les mêmes sept fichiers et les mêmes six
questions.
\end{nkbilan}

\begin{nkquestions}
\begin{enumerate}
\item Pourquoi écrit-on le banc avant le dessin, et qu'est-ce qui le rend
      possible ?
\item Quelles sont les deux seules inclusions du fichier de règles ?
\item Pourquoi une pièce capturée n'est-elle pas retirée du damier pendant
      l'exploration d'une rafle ?
\item Pourquoi le filtrage par rafle maximale se fait-il en sortie et non
      pendant l'exploration ?
\item Que se passe-t-il si l'on oublie le cas « camp sans coup légal » ?
\item Où s'exécute \texttt{OnCommandLine}, et pourquoi est-ce décisif ?
\end{enumerate}
\end{nkquestions}

\begin{nkpratique}
Construisez le \textbf{morpion} en suivant les treize étapes, dans l'ordre, sans
en sauter une.

Contraintes :
\begin{enumerate}
\item votre \texttt{NkMorpionRegles.h} n'inclut que \texttt{NkVector.h} et
      \texttt{NkTypes.h} ;
\item votre banc tourne \emph{avant} que le moindre pixel soit dessiné, et
      contient au moins deux cas négatifs --- on ne joue pas sur une case
      occupée, et une grille pleine sans alignement est une partie nulle ;
\item les trois modes (humain/IA, humain/humain, IA/IA) sortent d'un tableau de
      deux contrôleurs, sans un seul \texttt{if} de mode.
\end{enumerate}

Puis construisez-le pour le Web et ouvrez-le dans un navigateur.
\end{nkpratique}

\begin{nkdemo}
Prouvez que votre banc de morpion mesure vraiment quelque chose.

Cassez exprès la détection d'alignement --- ne testez que les lignes, pas les
colonnes. Lancez : il doit passer au \textbf{rouge}, sur le cas des colonnes,
avec un code de sortie non nul. Restaurez, relancez, notez le retour au vert.

\textbf{Recopiez les deux sorties} dans votre compte rendu : c'est la paire qui
prouve, pas le vert seul.

⚠️ Si la mutation \emph{ne casse rien}, ne concluez pas trop vite que le contrôle
manque : vérifiez d'abord que votre jeu d'essai pouvait exprimer l'écart. Une
grille où l'alignement gagnant est toujours une ligne ne peut pas voir un défaut
sur les colonnes.
\end{nkdemo}

\begin{nklogique}
Un joueur signale : « au ludo, quand je désactive deux sièges, le jeu se fige
après quelques tours ».

\begin{enumerate}
\item Sans lire le code, énumérez au moins trois causes possibles, en
      distinguant celles qui vivraient dans les \emph{règles} de celles qui
      vivraient dans l'\emph{écran}.
\item Pour chacune, dites quel cas de banc l'aurait attrapée.
\item Concluez sur celles qu'aucun banc de règles ne peut voir --- et dites par
      quel autre moyen on les prend.
\end{enumerate}
\end{nklogique}
